\section{\ComputationCircuits{}}

With our aim to sample from \CompActNets{}, let us first investigate mechanisms to represent the computation network.
More precisely, we construct sparse circuits preparing q-samples to the normalized computation network
\begin{align*}
    \probat{\headvariables,\shortcatvariables}
    = \frac{1}{\contraction{\bencodingofat{\sstat}{\headvariables,\shortcatvariables}}} \cdot \bencodingofat{\sstat}{\headvariables,\shortcatvariables} \, .
\end{align*}

\begin{definition}[\ComputationCircuit{}]
    Given a boolean function $\exformula:\atomstates\rightarrow[2]$, its oracle or \computationCircuit{} is the unitary tensor
    \begin{align*}
        \qcbencodingofat{\exformula}{\cheadvariableof{\exformula}{\insymbol},\cheadvariableof{\exformula}{\outsymbol},\shortcatvariables}
        &\quad=
        \sum_{\shortcatindicesin\wcols\exformulaat{\shortcatindices}=1}
        \paulixat{\cheadvariableof{\exformula}{\insymbol},\cheadvariableof{\exformula}{\outsymbol}} \otimes
        \onehotmapofat{\shortcatindices}{\shortcatvariables} \\
        & \quad \quad \quad +
        \sum_{\shortcatindicesin\wcols\exformulaat{\shortcatindices}=0}
        \identityat{\cheadvariableof{\exformula}{\insymbol},\cheadvariableof{\exformula}{\outsymbol}} \otimes
        \onehotmapofat{\shortcatindices}{\shortcatvariables} \, .
    \end{align*}
\end{definition}

\begin{example}[CNOT gate]\label{exa:cnotBenc}
    The CNOT gate is the controlled Pauli-X, that is
    \begin{align*}
        \mathrm{CNOT}\left[\cheadvariable{\insymbol},\cheadvariable{\outsymbol},\catvariable\right]
        = \paulixat{\cheadvariable{\insymbol},\cheadvariable{\outsymbol}} \otimes \tbasisat{\catvariable}
        + \identityat{\cheadvariable{\insymbol},\cheadvariable{\outsymbol}} \otimes \fbasisat{\catvariable} \, .
    \end{align*}
    For the identity function
    \begin{align*}
        \mathrm{Id}(x) = x \, .
    \end{align*}
    we have
    \begin{align*}
        \qcbencodingofat{\mathrm{Id}}{\cheadvariable{\insymbol},\cheadvariable{\outsymbol},\catvariable}
        = \mathrm{CNOT}\left[\cheadvariable{\insymbol},\cheadvariable{\outsymbol},\catvariable\right] \, .
    \end{align*}
    The $\mathrm{CNOT}$ is therefore the \computationCircuit{} to the identity function.
\end{example}

% Multiple output variables
Functions with multiple output variables, i.e. $\exformula:\atomstates\rightarrow\bigtimes_{\selindexin}[2]$, can be encoded image coordinate wise as a concatenation of the respective circuits.

\subsection{Relation with basis encodings}

\BasisEncodings{} are schemes to encode functions into directed boolean tensors, which enable reasoning about functions by tensor network contractions (see \cite{goessmann_tensor_2026}).
To connect with this formalism, let us now show the relation to the \computationCircuits{}.

\begin{lemma}
    \label{lem:bencComCircuit}
    For any boolean function we have (see \figref{fig:bencComCircuit})
    \begin{align*}
        \qcbencodingofat{\exformula}{\cheadvariableof{\exformula}{\insymbol},\cheadvariableof{\exformula}{\outsymbol},\shortcatvariables}
        = \contractionof{
            \bencodingofat{\exformula}{\headvariableof{\exformula},\shortcatvariables},
            \bencodingofat{\oplus}{\cheadvariableof{\exformula}{\outsymbol},\cheadvariableof{\exformula}{\insymbol},\headvariableof{\exformula}}
        }{\cheadvariableof{\exformula}{\insymbol},\cheadvariableof{\exformula}{\outsymbol},\shortcatvariables} \, .
    \end{align*}
\end{lemma}


% Interpretation as q-sample
The \computationCircuit{} applied on the initial state
\begin{align*}
    \sqrt{\frac{1}{2^{\catorder}}} \cdot \onesat{\ccatvariableof{[\atomorder]}{\outsymbol}} \otimes \fbasisat{\cheadvariable{\insymbol}}
\end{align*}
prepares a q-sample to the normalized basis encoding. % Called that basis encoding state!


% Sparse decomposition
In \cite{goessmann_tensor_2026} (and more detailed in \cite{goessmann_tensor-network_2025}) sparse decompositions of \basisEncodings{} have been studied.
Based on the equivalence of \computationCircuits{} and contracted \basisEncodings{}, stated by \lemref{lem:bencComCircuit} we transfer these decomposition schemes in the following.

\begin{figure}[t]
    \begin{center}
        \input{tikz_pics/computation-circuits/bencoding_computation_circuit}
    \end{center}
    \caption{Relation of computation circuit and basis encodings.
    The \computationCircuit{} of a function is the contraction of its \basisEncoding{} with the \basisEncoding{} of the mod2 sum $\oplus$.}
    \label{fig:bencComCircuit}
\end{figure}

% Good place?
\begin{example}[Multiple-controlled NOT]\label{exa:mcnotBenc}
    Let us consider the $\catorder$-ary controlled $\mathrm{NOT}$ gate
    \begin{align*}
        \mathrm{MCNOT}^{\catorder}\left[\cheadvariable{\insymbol},\cheadvariable{\outsymbol},\shortcatvariables\right]
        = \paulixat{\cheadvariable{\insymbol},\cheadvariable{\outsymbol}} \otimes \onehotmapofat{\onetuple{[\catorder]}}{\shortcatvariables}
        + \identityat{\cheadvariable{\insymbol},\cheadvariable{\outsymbol}} \otimes \left(\onesat{\shortcatvariables}-\onehotmapofat{\onetuple{[\catorder]}}{\shortcatvariables}\right) \, .
    \end{align*}
    This is the \computationCircuit{} of the $\catorder$-ary logical conjunction
    \begin{align*}
        \land^{\catorder} \defcols [2]^{\catorder} \rightarrow [2]
        \quad , \quad
         \land^{\catorder}(\shortcatindices)
        = \begin{cases}
              1 & \ifspace \shortcatindices = \onetuple{[\catorder]} \\
              0 & \text{else}
        \end{cases} \, .
    \end{align*}
    Special instances are the $\mathrm{CNOT}$ for $\catorder=1$ (see \exaref{exa:cnotBenc}) and the $\mathrm{Toffoli}$ gate for $\catorder=2$.
\end{example}

\subsection{Decomposition of functions}

The decomposition by contraction property of basis encodings is now a composition of circuits property, as stated in the next lemma.
Instead of hidden variables as in the \tnreason{} approach, we here exploit further working qubits.

\begin{lemma}
    \label{lem:comCirDecSpar}
    We have for functions $\exfunction:\atomstates\rightarrow\bigtimes_{\selindexin}[2]$, $\secexfunction:\bigtimes_{\selindexin}[2]\rightarrow\bigtimes_{s\in[r]}[2]$ (see \figref{fig:qcbencodingDecomposition})
    \begin{align*}
        &\qcbencodingofat{\secexfunction\circ\exfunction}{
            \cheadvariableof{\secexfunction\circ\exfunction}{\insymbol},
            \cheadvariableof{\secexfunction\circ\exfunction}{\outsymbol},
            \shortcatvariables}
        \otimes \fbasisat{\cheadvariableof{\exfunction}{\outsymbol}} \\
        &\quad= \breakablecontractionof{\onehotmapofat{0}{\cheadvariableof{\exfunction}{\insymbol}},
            \qcbencodingofat{\exfunction}{\cheadvariableof{\exfunction}{\insymbol},
                \cheadvariableof{\exfunction}{\midsymbol},
                \shortcatvariables},
            \qcbencodingofat{\secexfunction}{
                \headvariableof{\insymbol,\secexfunction\circ\exfunction},
                \cheadvariableof{\secexfunction\circ\exfunction}{\outsymbol},
                \cheadvariableof{\exfunction}{\midsymbol}},
            \qcbencodingofat{\exfunction}{\cheadvariableof{\exfunction}{\midsymbol},
                \cheadvariableof{\exfunction}{\outsymbol},
                \shortcatvariables}
        }{
            \cheadvariableof{\secexfunction\circ\exfunction}{\insymbol},\cheadvariableof{\secexfunction\circ\exfunction}{\outsymbol},\shortcatvariables,\cheadvariableof{\exfunction}{\outsymbol}
        } \, .
    \end{align*}
    Here each copy of $\headvariableof{\secexfunction\circ\exfunction}$ denotes a tuple of $r$ boolean variables and $\headvariableof{\exfunction}$ of $\seldim$ boolean variables.
\end{lemma}
\begin{proof}
    We show the claim based on basis calculus and sketch the steps in \figref{fig:proofDecomposition}.
    In the first equation we use \lemref{lem:bencComCircuit} and represent each \computationCircuit{} by a contraction of two basis encodings.
    In the second equation we use the identity
    \begin{align*}
        \contractionof{\bencodingofat{\oplus}{\cheadvariable{\outsymbol},\cheadvariable{\insymbol},\headvariableof{\exfunction}},
            \fbasisat{\cheadvariable{\insymbol}}}{\cheadvariable{\outsymbol},\headvariableof{\exfunction}}
        = \identityat{\cheadvariable{\outsymbol},\headvariableof{\exfunction}} \, .
    \end{align*}
    Now in the third equation we use that the contraction of basis encodings is the basis encoding of the composition function, and that $\exfunction\oplus\exfunction$ is the vanishing function, which basis encoding is a tensor product with $\fbasisat{\cheadvariableof{\exfunction}{\outsymbol}}$.
    We use we \lemref{lem:bencComCircuit} again in the fourth equation to arrive at the claim.
%    For any $\shortcatindices\in[2]^{\atomorder}$ we have that (see \figref{fig:proofDecomposition}):
%    \red{Uses the polynomial sparsity to uncompute: $\exfunction \oplus\exfunction = 0$.}
%    Can be seen by the basis encoding representation, when starting the hidden qubit with $\fbasis$, since
%    \begin{align*}
%        \contractionof{\bencodingofat{\oplus}{\headvariable,X_1},\fbasisat{X_0}}{Y,X_1}
%        = \identityat{Y,X_1} \, .
%    \end{align*}
%    The marginalization of the hidden target qubit $X_1$ for measurements does not change the distribution, since we have a deterministic dependence on the measured data register (need preparation by $\fbasis$ for that).
\end{proof}

When having a syntactical decomposition of a propositional formula, we can iteratively apply the \computationCircuit{} decomposition theorem and prepare each connective by a circuit.
We can decompose any propositional formula into logical connectives and prepare to each a \computationCircuit{} (exploiting further sparsity mechanisms).
This works, when the target qubit of one connective is used as a value qubit of another.

%% Closure of
This entanglement of the working qubit with the distributed qubits can be resolved by an un-computation circuit, which is the adjoint of the \computationCircuit{} for the working qubits (see \figref{fig:proofDecomposition}).
Note that while this is necessary in some applications, it is not necessary for the quantum rejection sampling presented here.

%\begin{figure}
%    \begin{center}
%        \input{tikz_pics/computation-circuits/decomposition_sparsity}
%    \end{center}
%    \caption{
%        Exploitation of \DecompositionSparsity{} in \computationCircuit{}s.
%        The working qubit $\headvariableof{\exfunction}$ is initialized in the ground state $\fbasis$ and un-computed after its usage to prepare $\headvariableof{\secexfunction\circ\exfunction}$.
%        It is then disentangled with the other qubits.
%    }\label{fig:qcbencodingDecomposition}
%\end{figure}

\begin{figure}
    \begin{center}
        \input{tikz_pics/computation-circuits/proof_decomposition_sparsity}
    \end{center}
    \caption{
        Proof of \lemref{lem:comCirDecSpar} based on basis calculus.
    }\label{fig:proofDecomposition}
\end{figure}

\begin{example}\label{exa:mcnotDecomposition}
    Let us reconsider the $\catorder$-ary controlled $\mathrm{NOT}$ gate from \exaref{exa:mcnotBenc}.
    The $\catorder$-ary logical conjunction is decomposed as a sequence of $2$-ary logical conjunctions.
    The \computationCircuit{} of each $2$-ary conjunction is a $\mathrm{Toffoli}$ gate applied to an input qubit and a worker qubit.
    Using \DecompositionSparsity{} we therefore find the standard decomposition of $\catorder$-ary controlled $\mathrm{NOT}$ gates into $\mathrm{Toffoli}$ gates.
\end{example}

\subsection{Reed-Muller Expansions}

% Comparison with tnreason
%This encoding strategy exploits a modified (by mod2 calculus) \polynomialSparsity{}.
We now investigate analogons to \polynomialSparsity{}.
First of all, we notice that the concatenation of \computationCircuits{} is a \computationCircuit{} to the mod2 sum.

\begin{lemma} %% Redundant: This follows from the oplus
    \label{lem:comCircuitCommutative}
    The composition of two \computationCircuits{} to functions $\exfunction$ and $\secexfunction$ (see \figref{fig:mod2ComputationCircuit}), where the \computationCircuits{} share the target qubit, is the \computationCircuit{} of their sum mod2.
    In particular, \computationCircuits{} commute when they have the same target qubit.
\end{lemma}
\begin{proof}
    This follows from the representation of \computationCircuits{} by basis encodings of the corresponding function and the mod2 sum $\oplus$ (see \figref{fig:mod2ComputationCircuit}).
\end{proof}

\begin{figure}[t]
    \begin{center}
        \input{tikz_pics/computation-circuits/mod2_sum_computation_circuit}
    \end{center}
    \caption{Equivalence of the composition of \computationCircuits{} with the \computationCircuit{} of their mod2 sum (\lemref{lem:comCircuitCommutative}).}
    \label{fig:mod2ComputationCircuit}
\end{figure}

% General application
When \computationCircuit{}s are available for the subformulas, one can therefore concatenatet them to construct the \computationCircuit{} of their mod2 sum.
Let us now suggest a set of basis functions, namely Mixed Polarity Reed-Muller (MPRM) terms to exploit this decomposition scheme.

\begin{definition}
    The MPRM term to $\hardparam$, where $\hardlegset\subset[\catorder]$ and $\headindexof{\hardlegset}\in\bigtimes_{\catenumerator\in\hardlegset}[2]$, is the function
    \begin{align*}
        \formulaof{\hardparam}(\shortcatindices)
        = \left(\prod_{\catenumerator\in\hardlegset\ncond\headindexof{\catenumerator}=1} \catindexof{\catenumerator}\right)
        \cdot \left(\prod_{\catenumerator\in\hardlegset\ncond\headindexof{\catenumerator}=0} (1-\catindexof{\catenumerator})\right) \, .
    \end{align*}
    A MPRM expansion of a formula $\formula$ is a set $\sliceset$ of tuples $\hardparam$ such that
    \begin{align*}
        \formula = \bigoplus_{\hardparam\in\sliceset} \formulaof{\hardparam} \, .
    \end{align*}
\end{definition}

% Algebraic Normal Form
This decomposition is closely related to the Algebraic Normal Form of a Boolean polynomial.
Let us now determine \computationCircuits{} for MPRM terms and therefore for arbitrary formulas in MPRM expansions.

\begin{lemma}
    A MPRM term $\formulaof{\hardparam}$ has the \computationCircuit{} (see \figref{fig:qcbencodingPolynomial})
    \begin{align*}
        &\contractionof{\qcaencodingofat{\formulaof{\hardparam}}{\cheadvariable{\insymbol},\cheadvariable{\outsymbol},\shortcatvariables},
            \identityat{\shortcatvariables,\ccatvariableof{[\catorder]}{\insymbol},\ccatvariableof{[\catorder]}{\outsymbol}}}{
            \cheadvariable{\insymbol},\cheadvariable{\outsymbol},\ccatvariableof{[\catorder]}{\insymbol},\ccatvariableof{[\catorder]}{\outsymbol}
        } \\
        &\quad = \breakablecontractionof{
            \{\mathrm{MCNOT}^{\cardof{\hardlegset}}[\cheadvariable{\insymbol},\cheadvariable{\outsymbol},\catvariableof{\hardlegset}]\}
            \cup \left(\bigcup_{\catenumerator\in\hardlegset\ncond\headindexof{\catenumerator}=0}\{\paulixat{\ccatvariableof{\catenumerator}{\insymbol},\catvariableof{\catenumerator}},
                     \paulixat{\catvariableof{\catenumerator},\ccatvariableof{\catenumerator}{\outsymbol}}\}\right) \\
            &\quad\quad\quad\quad\quad \cup \left(\bigcup_{\catenumerator\in\hardlegset\ncond\headindexof{\catenumerator}=1}\{
                                                \identityat{\catvariableof{\catenumerator},\ccatvariableof{\catenumerator}{\insymbol},\ccatvariableof{\catenumerator}{\outsymbol}}
                                                \}\right)
        }{
            \cheadvariable{\insymbol},\cheadvariable{\outsymbol},\ccatvariableof{[\catorder]}{\insymbol},\ccatvariableof{[\catorder]}{\outsymbol}
        } \, .
    \end{align*}
    The \computationCircuit{} of a formula $\formula$ with MPRM expansion $\sliceset$ is a concatenation of the \computationCircuits{} to $\hardlegset\in\sliceset$.
\end{lemma}
\begin{proof}
    For any $\ccatindexof{[\catorder]}{\insymbol},\ccatindexof{[\catorder]}{\outsymbol}$ the slices on both sides are vanishing, if $\ccatindexof{[\catorder]}{\insymbol}\neq\ccatindexof{[\catorder]}{\outsymbol}$.
    If $\ccatindexof{[\catorder]}{\insymbol}=\ccatindexof{[\catorder]}{\outsymbol}$ we have
    \begin{align*}
        &\breakablecontractionof{
            \{\mathrm{MCNOT}^{\cardof{\hardlegset}}[\cheadvariable{\insymbol},\cheadvariable{\outsymbol},\catvariableof{\hardlegset}]\}
            \cup \left(\bigcup_{\catenumerator\in\hardlegset\ncond\headindexof{\catenumerator}=0}\{\paulixat{\ccatvariableof{\catenumerator}{\insymbol},\catvariableof{\catenumerator}},
                     \paulixat{\catvariableof{\catenumerator},\ccatvariableof{\catenumerator}{\outsymbol}}\}\right) \\
            &\quad\quad\quad\quad \cup \left(\bigcup_{\catenumerator\in\hardlegset\ncond\headindexof{\catenumerator}=1}\{
                                           \identityat{\catvariableof{\catenumerator},\ccatvariableof{\catenumerator}{\insymbol},\ccatvariableof{\catenumerator}{\outsymbol}}
                                           \}\right)
        }{
            \cheadvariable{\insymbol},\cheadvariable{\outsymbol},\ccatvariableof{[\catorder]}{\insymbol}=\ccatindexof{[\catorder]}{\insymbol},\ccatvariableof{[\catorder]}{\outsymbol}=\ccatindexof{[\catorder]}{\outsymbol}
        } \\
        & \quad = \contractionof{\{\mathrm{MCNOT}^{\cardof{\hardlegset}}[\cheadvariable{\insymbol},\cheadvariable{\outsymbol},\catvariableof{\hardlegset}]\}
            \cup\{\onehotmapofat{(1-\ccatindexof{\catenumerator}{\insymbol})}{\catvariableof{\catenumerator}} \wcols \catenumerator\in\hardlegset\ncond\headindexof{\catenumerator}=0\}
            \cup\{\onehotmapofat{\ccatindexof{\catenumerator}{\insymbol}}{\catvariableof{\catenumerator}} \wcols \catenumerator\in\hardlegset\ncond\headindexof{\catenumerator}=1\}
        }{ \cheadvariable{\insymbol},\cheadvariable{\outsymbol}}\\
        & \quad =
        \begin{cases}
            \mathrm{NOT}[\cheadvariable{\insymbol},\cheadvariable{\outsymbol}] & \ifspace \formulaof{\hardparam}(\ccatindexof{[\catorder]}{\insymbol})=1 \\
            \identityat{\cheadvariable{\insymbol},\cheadvariable{\outsymbol}} & \ifspace \formulaof{\hardparam}(\ccatindexof{[\catorder]}{\insymbol})=0
        \end{cases} \, .
    \end{align*}
    Thus, the right contraction coincides for all slices on the variables $\ccatvariableof{[\catorder]}{\insymbol}$ and $\ccatvariableof{[\catorder]}{\outsymbol}$ with the oracle of $\formulaof{\hardparam}$.
\end{proof}

\begin{figure}
    \begin{center}
        \input{tikz_pics/computation-circuits/MPRM-term-cc}
    \end{center}
    \caption{
        \ComputationCircuit{} to a MPRM term $\formulaof{\hardparam}$, consisting in a multiple controlled $\mathrm{NOT}$ gate and Pauli-X gates for legs with $\headindexof{\catenumerator}=0$.
%        Representation of a slice as a controlled $\mathrm{CNOT}$ gate.
%        When the control state to a variable is $0$, we introduce auxiliary Pauli-X gates.
    }\label{fig:qcbencodingPolynomial}
\end{figure}


